\chapter{Instalasi}

\section{Metode 1: Install dari Go (Direkomendasikan)}

Cara termudah adalah menggunakan \cmd{go install} untuk mengunduh dan mengkompilasi
PakAI secara langsung:

\begin{lstlisting}[style=terminal]
$ go install github.com/dhanifudin/pakai/cmd/pakai@latest
\end{lstlisting}

Pastikan direktori \texttt{\$GOPATH/bin} (biasanya \texttt{\textasciitilde/go/bin})
sudah ada di \texttt{\$PATH}:

\begin{lstlisting}[style=terminal]
# Tambahkan ke ~/.bashrc atau ~/.zshrc
export PATH="$PATH:$HOME/go/bin"
\end{lstlisting}

\section{Metode 2: Build dari Kode Sumber}

\subsection{Kloning repositori}

\begin{lstlisting}[style=terminal]
$ git clone https://github.com/dhanifudin/pakai.git
$ cd pakai
\end{lstlisting}

\subsection{Kompilasi}

\begin{lstlisting}[style=terminal]
$ go build -ldflags="-s -w" -o pakai ./cmd/pakai
\end{lstlisting}

Atau gunakan Makefile yang sudah disediakan:

\begin{lstlisting}[style=terminal]
$ make build
\end{lstlisting}

\subsection{Pasang ke PATH}

\begin{lstlisting}[style=terminal]
$ sudo mv pakai /usr/local/bin/
# atau ke direktori lokal tanpa sudo:
$ mv pakai ~/.local/bin/
\end{lstlisting}

\section{Metode 3: Binary Release}

Binary siap pakai tersedia di halaman \textit{Releases} GitHub untuk arsitektur
Linux AMD64 dan ARM64:

\begin{enumerate}
  \item Kunjungi \url{https://github.com/dhanifudin/pakai/releases}
  \item Unduh arsip yang sesuai, misalnya:\\
        \texttt{pakai\_1.0.0\_linux\_amd64.tar.gz}
  \item Ekstrak dan pindahkan binary:
\end{enumerate}

\begin{lstlisting}[style=terminal]
$ tar -xzf pakai_1.0.0_linux_amd64.tar.gz
$ chmod +x pakai
$ mv pakai ~/.local/bin/
\end{lstlisting}

\section{Verifikasi Instalasi}

Setelah instalasi, mulai daemon lalu jalankan perintah status untuk memverifikasi
PakAI berjalan dengan benar:

\begin{lstlisting}[style=terminal]
$ pakai daemon start
$ pakai status
\end{lstlisting}

\begin{figure}[h]
  \centering
  \includegraphics[width=0.85\linewidth]{pakai-status}
  \caption{Output \texttt{pakai status} setelah instalasi berhasil}
\end{figure}

Jika perintah tidak ditemukan, pastikan direktori instalasi ada di \texttt{\$PATH}.

\begin{warningbox}
  Jika menggunakan \textbf{asdf} atau \textbf{mise} sebagai manajer versi Go,
  pastikan shim untuk \texttt{pakai} sudah diperbarui setelah instalasi.
  Jalankan \cmd{asdf reshim golang} atau \cmd{mise reshim} setelah proses instalasi.
\end{warningbox}
